\chapter{Técnica de prova FGV --- como atacar qualquer questão}

Guia transversal, válido para todas as matérias, destilado da resolução de
${\sim}$500 questões reais da FGV em 7 provas. Os capítulos por assunto dizem
\emph{o que} cada matéria cobra; este guia diz \emph{como} atacar a questão
na hora da prova. Leia-o primeiro e volte a ele na última semana.

\section{Leia o comando antes de tudo}

O comando é a última linha do enunciado. O erro número um em prova é responder
o oposto do que foi pedido.

\begin{itemize}
\item \textbf{``Assinale a correta''} $\times$ \textbf{``assinale a
INCORRETA / a que NÃO...''} --- a FGV usa muito o comando negativo. Circule o
``incorreta''/``não'' no papel antes de ler as alternativas.
\item \textbf{``Julgue as afirmativas I, II, III''} $\to$ resolva cada item
isolado como V/F, depois case com a alternativa (``I e II apenas'', ``I, II e
III'', etc.). Uma afirmativa falsa derruba toda alternativa que a inclua ---
use isso para eliminar em bloco.
\item \textbf{Sequência (V) (F) (F)}: mesma lógica. Acertar 1 ou 2 itens já
elimina metade das alternativas.
\end{itemize}

\section{Os padrões de distrator da FGV (o coração do método)}

A FGV constrói a alternativa errada de poucas formas. Reconhecê-las é meio
caminho andado.

\begin{pegadinha}[Os sete padrões de distrator]
\begin{enumerate}
\item \textbf{Absolutos.} ``sempre'', ``nunca'', ``exclusivamente'',
``apenas'', ``todo'', ``invariavelmente'', ``impossível'', ``garante''. No
mundo técnico quase nada é absoluto --- a alternativa com absoluto é quase
sempre o distrator.
\item \textbf{Inversão de conceitos.} Troca dois conceitos do mesmo par:
OLTP$\leftrightarrow$OLAP, PUT$\leftrightarrow$PATCH,
agregação$\leftrightarrow$composição, checked$\leftrightarrow$unchecked,
controlador$\leftrightarrow$operador, governança$\leftrightarrow$gestão,
IDS$\leftrightarrow$IPS, SAST$\leftrightarrow$DAST, TCP$\leftrightarrow$UDP,
incidente$\leftrightarrow$problema, Factory$\leftrightarrow$Abstract Factory.
Se você domina o par lado a lado, a inversão salta aos olhos.
\item \textbf{A ``quase certa''.} Acerta a primeira metade e erra a segunda
(``X é um índice bitmap, adequado para \emph{alta} cardinalidade'' --- bitmap
é para \emph{baixa}). Leia a alternativa inteira; a FGV esconde o erro no
fim.
\item \textbf{Extrapolação (interpretação).} Em português/inglês, o distrator
diz mais do que o texto diz, ou inverte a tese do autor. Se não está
sustentado no texto, é distrator --- não importa se é verdade no mundo real.
\item \textbf{Troca de número.} ``5 domínios do COBIT'', ``4 dimensões do
ITIL'', ``3 formas normais'', ``camada 7 do OSI''. A FGV troca o número.
Decore as quantidades.
\item \textbf{Troca de ordem.} Fases do ARIES (Analysis $\to$ Redo $\to$
Undo), ciclo de vida, camadas, formas normais. O distrator embaralha a
sequência.
\item \textbf{Contradição interna.} ``SOAP sem contrato formal'' (SOAP usa
WSDL), ``microsserviço com banco único compartilhado'' (contra o princípio).
A alternativa se contradiz --- descarte.
\end{enumerate}
\end{pegadinha}

\section{Estratégia de resolução}

\begin{enumerate}
\item \textbf{Elimine primeiro os absolutos e as contradições internas} ---
costumam sumir 2 alternativas de cara.
\item \textbf{Compare as 2 que sobraram.} A FGV quase sempre deixa uma
``quase certa'' como pegadinha final; a diferença está num detalhe (uma
palavra, a segunda metade da frase, um número).
\item \textbf{Volte ao enunciado} e confirme a escolhida contra o que foi
pedido (o cenário e o comando), não contra sua memória solta.
\item \textbf{Interpretação:} volte ao trecho exato. Nunca responda ``pelo
que faz sentido''; responda pelo que o texto afirma.
\item \textbf{Só no chute, empatado: fique com a mais longa.} Medido nas 15
provas FGV do banco de questões: a correta é a alternativa mais longa em
\textbf{33\%} dos itens, contra 20\% de acaso --- e o efeito se concentra nos
blocos técnicos (BI 50\%, banco de dados 49\%, programação 46\%, atualidades
41\%). Em português (25\%) e inglês (17\%) \textbf{não vale}: ali a banca
nivela o tamanho das alternativas. É desempate de último recurso, depois de
esgotar conteúdo e eliminação --- nunca critério de escolha.
\end{enumerate}

\section{Gestão de tempo (70 questões)}

\textbf{Atenção:} o edital \textbf{não declara a duração da prova}. O que ele
fixa é o fechamento dos portões (12h30), a \textbf{permanência mínima de 2
horas} (item 10.7) e a entrega do caderno só nos últimos 30 minutos (10.9). O
padrão da FGV para 70 questões é \textbf{4 horas}, e é com esse número que o
ritmo abaixo foi calculado --- mas \textbf{confirme no edital de convocação}
antes de calibrar o treino.

\begin{comosair}[Ritmo de prova]
\begin{itemize}
\item \textbf{Específicos primeiro.} Valem 2,5$\times$ (75 dos 115 pontos).
Comece por eles com a cabeça fresca; deixe português/inglês/RLM para depois.
\item \textbf{Não deixe disciplina em branco.} Zerar qualquer uma elimina,
independentemente da nota total (edital, 9.17 ``b'') --- reserve os últimos
minutos para não deixar nenhum bloco sem resposta.
\item \textbf{Não trave.} Marque a questão difícil, siga, volte no fim. Uma
questão de RLM não vale mais que uma de eng.\ de software --- e costuma
custar mais tempo.
\item \textbf{Ritmo de referência:} ${\sim}$3 min por questão em média
(supondo 4h). Treine com \texttt{./quiz.py --simulado} (cronometrado, sem
feedback até o fim).
\end{itemize}
\end{comosair}

\section{Regras de ouro}

\begin{itemize}
\item \textbf{Filtro travado em FGV.} Cada banca tem uma gramática de
pegadinha própria; reflexo de CESPE ou de outra banca atrapalha aqui.
\item \textbf{Cartão de resposta:} questão com duas marcações ou nenhuma =
zero. Confira. Reserve tempo para preencher com calma.
\item \textbf{Não mude resposta no chute.} Só troque se tiver um motivo
concreto (releu e viu o erro), não por insegurança de última hora.
\item \textbf{Anulada acontece.} Se a questão parecer ter duas certas ou
nenhuma, marque a ``menos errada'' e siga --- não perca tempo brigando com
ela.
\end{itemize}

\section{Depois de cada simulado}

Toda questão errada vira material de revisão: \texttt{./quiz.py --erradas}
(repetição espaçada) e a anotação automática em \texttt{erros/<bloco>.md}. O
caderno de erros é o material principal de revisão nas duas últimas semanas
antes da prova.
